\section{Background: pairs, products, and 2-safety}
\label{part:background}
\label{sec:twosafety}

This part builds the one piece of machinery the rest of the document leans on.
If you are fluent in information-flow policy but have not thought much about
\emph{how} a relational property gets checked mechanically, this is the gap
worth closing before Part~\ref{part:walkthrough}.

\subsection{Safety properties and their limits}

Most program verification concerns \emph{safety} properties: a predicate over a
single execution, of the form ``no execution ever reaches a bad state.'' Null
dereference, array bounds, assertion violation, and functional correctness
against a specification are all of this shape. The reason this class is
comfortable is that it is checkable by reachability --- build an abstraction of
the state space, show the bad set is unreachable, done. Every mature tool we
have, from abstract interpreters to model checkers to SMT-backed verifiers, is
built for it.

Equation~\eqref{eq:ni} is not of this shape. No single execution of $P$ is a
counterexample to it. A violation is a \emph{pair}: two executions that agree on
everything the observer is entitled to know, and disagree on some observable.
This is the defining feature of a \emph{2-safety} property, also called a
hyperproperty of arity two.

\begin{remark}[Read \eqref{eq:ni} as a specification, not as an encoding]
\label{rem:notencoding}
Equation~\eqref{eq:ni} is stated as a whole-run implication, and it is correct as
an informal statement of intent. It is \textbf{not} how the checker encodes the
question, and turning it directly into a solver query is a known unsoundness.
\cref{sec:ledger} explains why and gives the construction that replaces it.
\end{remark}

\begin{remark}
The practical consequence is worth stating plainly: you cannot instrument $P$,
run it, and check an assertion. There is no assertion over a single run whose
failure means ``leak.'' Dynamic taint tracking, which is what LLVM's
DataFlowSanitizer does, therefore cannot decide \eqref{eq:ni}; it approximates
it by tracking labels through one execution and reporting when a labelled value
reaches a sink. That approximation has no declassification and no notion of
timing, which is why it is the wrong model here even though its label plumbing
is instructive.
\end{remark}

\subsection{The product construction}

There is a standard and very useful move: \textbf{2-safety of $P$ is ordinary
safety of a program that runs two copies of $P$ at once.} Build $P \times P$,
the \emph{product} or \emph{self-composition}, whose state is a pair of $P$
states, one per \emph{lane}. Couple the initial states on what the observer can
see \emph{at entry}, let the secrets range freely, and add a bad-state predicate
that fires when the lanes' observations diverge.

Precisely which facts couple the initial states matters a great deal, and is the
subject of \cref{sec:ledger}. For now note only what they are \emph{not}: they do
not include equality of any release the program has not performed yet.

\begin{center}
\begin{tikzpicture}[
  lane/.style={draw, rounded corners=1.5pt, minimum height=5.5mm, minimum width=14mm,
               font=\small, inner sep=2pt},
  bad/.style={draw, rounded corners=1.5pt, minimum height=5.5mm, minimum width=14mm,
              font=\small, fill=red!12, inner sep=2pt},
  lbl/.style={font=\scriptsize, anchor=center},
  ar/.style={-{Latex[length=1.6mm]}}
]
  \node[lane] (a0) at (0,0.8)   {$\sigma_0$};
  \node[lane] (a1) at (2.2,0.8) {$\sigma_0'$};
  \node[lane] (a2) at (4.4,0.8) {$\sigma_0''$};
  \node[bad]  (a3) at (6.7,0.8) {\textsf{obs}${}=\alpha$};

  \node[lane] (b0) at (0,-0.8)   {$\sigma_1$};
  \node[lane] (b1) at (2.2,-0.8) {$\sigma_1'$};
  \node[lane] (b2) at (4.4,-0.8) {$\sigma_1''$};
  \node[bad]  (b3) at (6.7,-0.8) {\textsf{obs}${}=\beta$};

  \foreach \i/\j in {a0/a1, a1/a2, a2/a3, b0/b1, b1/b2, b2/b3}
    \draw[ar] (\i) -- (\j);

  \foreach \i/\j in {a0/b0, a1/b1, a2/b2}
    \draw[dotted] (\i) -- (\j);

  \node[lbl, anchor=east] at (-0.85,0.8)  {lane 0: $s_0$};
  \node[lbl, anchor=east] at (-0.85,-0.8) {lane 1: $s_1$};

  \node[lbl] at (0,1.7)  {coupled at entry only};
  \node[lbl, text width=2.6cm, align=center] at (6.7,1.85)
       {$\alpha \neq \beta$ is the bad state};
\end{tikzpicture}
\captionof{figure}{The two-lane product. Initial states are coupled on what is
visible \emph{at entry} --- public configuration, public alias topology, and the
initially-visible components --- with the secrets free. Releases are \emph{not}
part of this coupling; they are handled step by step as the lanes run
(\cref{sec:ledger}). Dotted lines are the lockstep alignment. A violation is
exactly a reachable state where the lanes' projected observations differ, which
is an ordinary reachability question over the product.}
\label{fig:product}
\end{center}

This is the hinge of the design. It converts a property no existing tool can
check directly into one that every existing tool can check, at the cost of
squaring the state space. It is also why the specification speaks of an
``audit-all whole-entry product'' rather than a type system: the product is the
\emph{authoritative} object, and a verdict is a reachability result over it.

\subsection{The release ledger, and why whole-run equality is wrong}
\label{sec:ledger}

\begin{remark}[Two things called a ledger]
This document has a \emph{claim ledger} in the margin --- a presentational device
of these notes. The specification separately has a \emph{release ledger}, an
object inside the product semantics. They are unrelated. This subsection is about
the second; every margin panel is the first.
\end{remark}

\cref{rem:notencoding} promised an explanation. Here it is, because this is the
one place where the obvious encoding of \eqref{eq:ni} is not merely imprecise but
\emph{unsound}.

\textbf{What couples the initial states.} The lanes must agree on what the
coalition can see \emph{at entry}: the public configuration, the public alias
topology, and the read-carrier of every initially-visible component. Mechanism
and timing-environment choices are coupled too, so a pair cannot differ on a
modelled implementation choice. Deliberately absent: any requirement that the
secret components differ, and --- the point of this subsection --- any equality of
a release the program has not performed yet.

\textbf{Why that last omission matters.} \eqref{eq:ni} reads: if the authorized
releases come out equal, the observations must come out equal. The obvious
encoding installs $\Rel(p,s_0) = \Rel(p,s_1)$ as an initial constraint on the
product and then asks for bad-state reachability. A two-instruction program shows
that this is wrong:

\begin{center}
\code{emit(public\_channel, h)} \quad;\quad \code{release(policy, audience, h)}
\end{center}

The secret is published at step~1 and only afterwards passed through an authorized
release at step~2. Under the whole-run encoding the constraint keeps only lane
pairs whose complete release histories agree --- and since the released value
\emph{is} $h$, that means only pairs with equal $h$. Those pairs trivially agree at
step~1 too. The leak is never examined: the encoding has filtered away precisely
the pairs that demonstrate it. Note the direction of the error --- it is
permissive, not conservative.

\textbf{The fix is causal.} Releases are consulted step by step as the lanes run,
never installed up front:

\begin{center}
\begin{tabular}{@{}p{0.31\linewidth} p{0.58\linewidth}@{}}
\toprule
At an aligned step & The ledger does \\
\midrule
no release here
  & require the projected words to agree, as usual \\[3pt]
release authorized for this observer, values equal
  & append the entry and continue; the obligation stays active \\[3pt]
release authorized for this observer, values differ
  & permitted --- but everything outside the release's declared footprint must
    still agree, and the pair's obligation then \emph{retires} \\[3pt]
release not authorized for this observer
  & conceal the value; require everything else about the event to agree \\
\bottomrule
\end{tabular}
\end{center}

The transition takes only the current ledger and the current pair of events. It
has no parameter through which a later event can be inspected, so a release that
has not happened cannot excuse an observation that already has.

\begin{remark}
\label{rem:mtcm3}
``Retires'' is doing real work: an authorized difference ends the obligation for
that pair \emph{from that point}, rather than licensing differences
retroactively. A pair is protected up to its first authorized divergence and
released afterwards, which is what makes declassification expressible without
making it a loophole.

This is not hypothetical. The program above is encoded in the harness as
\fx{prefix\_causal\_release.bad.mlir}, and a checker built on the whole-run
encoding reports that fixture safe.
\end{remark}

\subsection{Why the product is harder than it looks}

Figure~\ref{fig:product} is drawn with the lanes in lockstep, which quietly
assumes they follow the same control path. They need not. Once a secret can
influence control flow the lanes desynchronize, and the product must say
something about misaligned executions. This is where the bad-state predicate
stops being ``the outputs differ'' and becomes a disjunction over structural
divergences:

\begin{enumerate}[nosep,leftmargin=1.6em]
\item one lane can step and the other cannot;
\item the two lanes' next control locations differ;
\item their termination or failure statuses differ;
\item their event site or occurrence alignment differs;
\item the projected event words they emit differ.
\end{enumerate}

Items 1--4 are \emph{structural}: they compare the shape of the two executions,
not their data. That distinction matters because structural divergence is not
weakened when the observer's projection conceals a payload. An observer who
cannot read a value can still count operations, notice that a branch was taken,
and see that a run terminated early.

\begin{remark}[Where the subtle bugs live]
Nearly every unsoundness in this area is a case where somebody checked item 5
and forgot items 1--4, or checked a \emph{local} version of item 4 rather than a
global one. Two of the countermodels in Part~\ref{part:walkthrough} are exactly
of this form, and one of them --- \cref{step:product} --- is a program where
every per-site payload and count agrees and only the global ordering differs.
\end{remark}

\subsection{Two layers, and the discipline between them}

Squaring the state space is expensive, so it is tempting to avoid the product
with a cheap syntactic analysis: label each value \textsf{Low} or \textsf{High},
join labels at merges, and complain when \textsf{High} reaches a sink. That is a
classic information-flow type system, and it is genuinely useful. But note
carefully which way its guarantee points. For a \emph{sound over-approximation}
it is \textbf{acceptance} that is conclusive --- the classical theorem is that a
well-typed program \emph{is} noninterferent --- while its \emph{rejections} may
be false alarms. A type system proves; it cannot refute.

Two things nevertheless remove that proving power in this setting. The
observation model here includes timing, addresses, allocation sizes, and
operation counts, and a label discipline over values says nothing about any of
them. And the property is release-relative, so a declassification must be
\emph{permitted} at exactly the right point, which the prefix-causal ledger of
\cref{sec:ledger} decides and a label lattice cannot express. What survives is a
useful \emph{triage}: cheap, sound in the direction of refusal, and unable to
close the question.

The architecture therefore has two layers with a deliberately asymmetric
contract:

\begin{center}
\begin{tabular}{@{}p{0.29\linewidth}p{0.61\linewidth}@{}}
\toprule
\textbf{Diagnostic layer} (\Lv{1}) & one \textsf{Low}/\textsf{High} pass, with
no proof-authoritative strong update, no function summaries, and no slice
selection. Deliberately weak. \\[3pt]
\textbf{Product layer} (\Lv{2}) & the exact two-lane product over the whole
entry. Authoritative. \\
\bottomrule
\end{tabular}
\end{center}

The contract is: \textbf{the diagnostic layer may narrow the work, and may never
discharge an obligation.} Its weakness is the point. Because it cannot perform a
strong update it cannot conclude that a secret store was later overwritten;
because it has no summaries it cannot conclude anything about a callee. So it
cannot be unsound --- it can only be imprecise, and its correct output at a site
it cannot decide is \RelReq{}: \emph{the product must decide this one}.

\begin{remark}[The single most dangerous repair]
\label{rem:danger}
When the diagnostic layer reports a violation on a program that is in fact safe,
the tempting fix is to teach it the missing reasoning --- a strong update, a
summary, a slice. Every such addition moves proof burden into the layer designed
to carry none, and the specification forbids the consequence outright: no
diagnostic disposition, including \emph{statically discharged} or \emph{not
observable}, may establish safety or remove a product constraint. This is why
the harness ships \emph{paired anti-controls} rather than controls alone, a
discipline developed in \cref{sec:pairs}.
\end{remark}
